Separate computations already address the earlier $n=1$, $l=0$ mode
near $\omega^2=0.7976767871$ and $\rho=1.7446175448$. They concern the
same sextic field equation, not C1 or C2. A finite-amplitude radial
evolution compared the nonlinear field with its background and linear
response on two spatial grids. For amplitudes $0.01,0.02,0.04$, the
integrated quadratic flux of the nonlinear remainder increased by
factors approximately $15.97$ and $15.89$ upon doubling amplitude.
This supports fourth-order scaling over the tested range; that remainder
flux is not an exact measurement of the total lost core energy.

A subsequent harmonic diagnostic identified the dominant laboratory
sideband $E=\omega-2\rho\simeq-2.59611$ as outgoing at two detector
radii, on both grids and in both prescribed time windows. The weaker
$E=\omega+2\rho$ component was resolved but did not satisfy every
outgoing-wave criterion. Large residuals in the full field's harmonic
fit remained, so these observations do not establish a periodic
multifrequency solution.

A separately solved frequency-domain problem, using the same background
and mode data but no fitted amplitude from the time evolution, predicted
the dominant outgoing coefficient. At $\epsilon=0.02$, its predicted
amplitude was $5.1563\times10^{-5}$, compared with measured values
approximately $(5.03$--$5.10)\times10^{-5}$. The prescribed comparison
criteria passed. Shared input profiles make this a separate mechanism
test, not a wholly independent reconstruction of the background.

An interval calculation further separated one necessary overlap from
zero for an explicitly frozen finite operator on $[0,44]$:
\begin{equation}
 I_2\in[-0.16383470526432667,-0.10547329373338812].
 \label{nl:frozen}
\end{equation}
The corresponding $I_1$ interval contained zero, and the original test
requiring both exclusions remained failed. Nevertheless, one certified
nonzero overlap suffices to violate the necessary simultaneous condition
$I_1=I_2=0$ for that frozen problem. Its binary64 spline inputs are treated
as exact data, with a free exterior and zero source beyond radius $44$.
This is not a certificate that the exact continuum Q-ball and its BIC
have that overlap: a validated link to the continuum inputs and exterior
tail remains required. The linear existence certificates do not, by
themselves, supply that quantitative link.

If a leading second-order outgoing amplitude is nonzero, its power is
of order $\epsilon^4$. A lifetime law additionally requires a controlled
relation to mode energy and slow amplitude evolution, and separation of
initial adjustment from persistent radiation. Neither the finite-time
observations nor the frozen-operator certificate provides a general
nonlinear stability or lifetime theorem. Earlier real-field shell-model
normal forms describe a different model and supply no substitute for
these Q-ball source and overlap conditions.
